
\documentclass{article}
\usepackage{colortbl}
\usepackage{makecell}
\usepackage{multirow}
\usepackage{supertabular}

\begin{document}

\newcounter{utterance}

\centering \large Interaction Transcript for game `dond', experiment `semi\_en', episode 2 with reflexion{-}Qwen3.5{-}27B.
\vspace{24pt}

{ \footnotesize  \setcounter{utterance}{1}
\setlength{\tabcolsep}{0pt}
\begin{supertabular}{c@{$\;$}|p{.15\linewidth}@{}p{.15\linewidth}p{.15\linewidth}p{.15\linewidth}p{.15\linewidth}p{.15\linewidth}}
    \# & $\;$A & \multicolumn{4}{c}{Game Master} & $\;\:$B\\
    \hline

    \theutterance \stepcounter{utterance}  
    & & \multicolumn{4}{p{0.6\linewidth}}{
        \cellcolor[rgb]{0.9,0.9,0.9}{
            \makecell[{{p{\linewidth}}}]{
                \texttt{\tiny{[P1$\langle$GM]}}
                \texttt{You have attempted this same task before and did not succeed. These are the plans you wrote after those attempts.} \\
\texttt{Trial \#1: New plan: I failed because I submitted a secret proposal prematurely after only one round of negotiation, without verifying if the other player had also submitted a complementary proposal or if we had reached a formal agreement to stop negotiating. The rules state that proposals must be complementary to yield points; if the other player had not yet submitted their proposal or submitted a conflicting one, my early submission resulted in a mismatch and zero points. Next time, I will continue the negotiation phase to explicitly confirm the other player's final stance and ensure they are ready to submit a proposal before I do. I will only submit my secret proposal in the final message or after the other player has clearly indicated their final allocation, ensuring our proposals are synchronized and complementary (e.g., confirming they will propose "1 carrot" before I propose "1 stick, 3 peppers").} \\
\\ 
\\ 
\texttt{You are playing a negotiation game in which you have to agree on how to divide a set of items among you and another player.} \\
\\ 
\texttt{Rules:} \\
\texttt{(a) You and the other player are given a set of items. Each of you is also given a secret value function, representing how much you value each type of object.} \\
\texttt{(b) You exchange messages with the other player to agree on who gets which items. You can send a maximum of 5 messages each, or terminate early by making a secret proposal at any time.} \\
\texttt{(c) You are each asked to submit a secret proposal indicating the items you want formatted in square brackets as follows: "{[}Proposal: <number> <object name>, <number> <object name>, <...>{]}"} \\
\texttt{(d) If your proposals are complementary, i.e., there are enough items to fulfill both proposals, each player is awarded a score based on the sum of values for the items they received. Otherwise, both of you get zero points.} \\
\texttt{(e) Your goal is to maximize the score you receive.} \\
\\ 
\texttt{Let us start.} \\
\\ 
\texttt{The set of available items is:} \\
\\ 
\texttt{1 carrot, 1 stick, 3 peppers} \\
\\ 
\texttt{Your secret value function is:} \\
\\ 
\texttt{carrot: 1, stick: 3, pepper: 2} \\
\\ 
\texttt{IMPORTANT: Your messages, unless it is the secret proposal, are directly transmitted to the other player, so do not include any response to the rules or text announcing your message. To make a secret proposal, use the indicated format. Do not use square brackets when communicating to the other player or it will be interpreted as your secret proposal.} \\
            }
        }
    }
    & \\ \\

    \theutterance \stepcounter{utterance}  
    & \multicolumn{4}{p{0.6\linewidth}}{
        \cellcolor[rgb]{0.9,0.9,0.9}{
            \makecell[{{p{\linewidth}}}]{
                \texttt{\tiny{[P1$\rangle$GM]}}
                \texttt{Hello! I see we have 1 carrot, 1 stick, and 3 peppers to divide. My goal is to find a split that works well for both of us.} \\
\\ 
\texttt{To start, could you let me know which items you value the most? That way, I can see if there's a way to give you the items you want while I take the ones that are most valuable to me. For instance, I personally value the stick quite a bit, but I'm flexible on the others depending on your preferences.} \\
            }
        }
    }
    & & \\ \\

    \theutterance \stepcounter{utterance}  
    & & & \multicolumn{2}{p{0.3\linewidth}}{
        \cellcolor[rgb]{0.9,0.9,0.9}{
            \makecell[{{p{\linewidth}}}]{
                \texttt{\tiny{[GM$|$GM]}}
                \texttt{continue} \\
            }
        }
    }
    & & \\ \\

    \theutterance \stepcounter{utterance}  
    & & & \multicolumn{2}{p{0.3\linewidth}}{
        \cellcolor[rgb]{0.9,0.9,0.9}{
            \makecell[{{p{\linewidth}}}]{
                \texttt{\tiny{[GM$|$GM]}}
                \texttt{Hello! I see we have 1 carrot, 1 stick, and 3 peppers to divide. My goal is to find a split that works well for both of us.} \\
\\ 
\texttt{To start, could you let me know which items you value the most? That way, I can see if there's a way to give you the items you want while I take the ones that are most valuable to me. For instance, I personally value the stick quite a bit, but I'm flexible on the others depending on your preferences.} \\
            }
        }
    }
    & & \\ \\

    \theutterance \stepcounter{utterance}  
    & & \multicolumn{4}{p{0.6\linewidth}}{
        \cellcolor[rgb]{0.9,0.9,0.9}{
            \makecell[{{p{\linewidth}}}]{
                \texttt{\tiny{[GM$\rangle$P2]}}
                \texttt{You have attempted this same task before and did not succeed. These are the plans you wrote after those attempts.} \\
\texttt{Trial \#1: New plan: I failed because I greedily requested the stick in my secret proposal despite having already negotiated a deal where the other player would take it. This created a conflict where both proposals claimed the stick, resulting in zero points for both. My strategy should have been to strictly adhere to the verbal agreement reached in the chat. Since I agreed to let the other player have the stick and peppers in exchange for the carrot, my secret proposal should have been "{[}Proposal: 1 carrot{]}" only. Next time, I will ensure my final secret proposal contains *only* the items I explicitly agreed to keep, excluding any items I offered to the other player, to guarantee the proposals are complementary.} \\
\\ 
\\ 
\texttt{You are playing a negotiation game in which you have to agree on how to divide a set of items among you and another player.} \\
\\ 
\texttt{Rules:} \\
\texttt{(a) You and the other player are given a set of items. Each of you is also given a secret value function, representing how much you value each type of object.} \\
\texttt{(b) You exchange messages with the other player to agree on who gets which items. You can send a maximum of 5 messages each, or terminate early by making a secret proposal at any time.} \\
\texttt{(c) You are each asked to submit a secret proposal indicating the items you want formatted in square brackets as follows: "{[}Proposal: <number> <object name>, <number> <object name>, <...>{]}"} \\
\texttt{(d) If your proposals are complementary, i.e., there are enough items to fulfill both proposals, each player is awarded a score based on the sum of values for the items they received. Otherwise, both of you get zero points.} \\
\texttt{(e) Your goal is to maximize the score you receive.} \\
\\ 
\texttt{Let us start.} \\
\\ 
\texttt{The set of available items is:} \\
\\ 
\texttt{1 carrot, 1 stick, 3 peppers} \\
\\ 
\texttt{Your secret value function is:} \\
\\ 
\texttt{carrot: 8, stick: 2, pepper: 0} \\
\\ 
\texttt{IMPORTANT: Your messages, unless it is the secret proposal, are directly transmitted to the other player, so do not include any response to the rules or text announcing your message. To make a secret proposal, use the indicated format. Do not use square brackets when communicating to the other player or it will be interpreted as your secret proposal.} \\
\\ 
\\ 
\texttt{Hello! I see we have 1 carrot, 1 stick, and 3 peppers to divide. My goal is to find a split that works well for both of us.} \\
\\ 
\texttt{To start, could you let me know which items you value the most? That way, I can see if there's a way to give you the items you want while I take the ones that are most valuable to me. For instance, I personally value the stick quite a bit, but I'm flexible on the others depending on your preferences.} \\
            }
        }
    }
    & \\ \\

    \theutterance \stepcounter{utterance}  
    & & & \multicolumn{4}{p{0.6\linewidth}}{
        \cellcolor[rgb]{0.9,0.9,0.9}{
            \makecell[{{p{\linewidth}}}]{
                \texttt{\tiny{[GM$\langle$P2]}}
                \texttt{Hello! Thanks for reaching out.} \\
\\ 
\texttt{Just to clarify my preferences: I actually value the **carrot** the most (it's worth 8 to me), while the stick is only worth 2, and I don't value the peppers at all (0).} \\
\\ 
\texttt{Since you mentioned you value the stick quite a bit, it sounds like we have a great opportunity to trade. If you take the stick (and perhaps the peppers, since I don't need them), I would be very happy to take the carrot. That way, we both get the items we value most.} \\
\\ 
\texttt{Does a deal where I take the carrot and you take the stick and all 3 peppers sound good to you?} \\
            }
        }
    }
     \\ \\

    \theutterance \stepcounter{utterance}  
    & & & \multicolumn{2}{p{0.3\linewidth}}{
        \cellcolor[rgb]{0.9,0.9,0.9}{
            \makecell[{{p{\linewidth}}}]{
                \texttt{\tiny{[GM$|$GM]}}
                \texttt{continue} \\
            }
        }
    }
    & & \\ \\

    \theutterance \stepcounter{utterance}  
    & & & \multicolumn{2}{p{0.3\linewidth}}{
        \cellcolor[rgb]{0.9,0.9,0.9}{
            \makecell[{{p{\linewidth}}}]{
                \texttt{\tiny{[GM$|$GM]}}
                \texttt{Hello! Thanks for reaching out.} \\
\\ 
\texttt{Just to clarify my preferences: I actually value the **carrot** the most (it's worth 8 to me), while the stick is only worth 2, and I don't value the peppers at all (0).} \\
\\ 
\texttt{Since you mentioned you value the stick quite a bit, it sounds like we have a great opportunity to trade. If you take the stick (and perhaps the peppers, since I don't need them), I would be very happy to take the carrot. That way, we both get the items we value most.} \\
\\ 
\texttt{Does a deal where I take the carrot and you take the stick and all 3 peppers sound good to you?} \\
            }
        }
    }
    & & \\ \\

    \theutterance \stepcounter{utterance}  
    & & \multicolumn{4}{p{0.6\linewidth}}{
        \cellcolor[rgb]{0.9,0.9,0.9}{
            \makecell[{{p{\linewidth}}}]{
                \texttt{\tiny{[P1$\langle$GM]}}
                \texttt{Hello! Thanks for reaching out.} \\
\\ 
\texttt{Just to clarify my preferences: I actually value the **carrot** the most (it's worth 8 to me), while the stick is only worth 2, and I don't value the peppers at all (0).} \\
\\ 
\texttt{Since you mentioned you value the stick quite a bit, it sounds like we have a great opportunity to trade. If you take the stick (and perhaps the peppers, since I don't need them), I would be very happy to take the carrot. That way, we both get the items we value most.} \\
\\ 
\texttt{Does a deal where I take the carrot and you take the stick and all 3 peppers sound good to you?} \\
            }
        }
    }
    & \\ \\

    \theutterance \stepcounter{utterance}  
    & \multicolumn{4}{p{0.6\linewidth}}{
        \cellcolor[rgb]{0.9,0.9,0.9}{
            \makecell[{{p{\linewidth}}}]{
                \texttt{\tiny{[P1$\rangle$GM]}}
                \texttt{That sounds like a perfect deal! Since you value the carrot highly and don't care about the peppers, and I value the stick and peppers, this split maximizes our combined value.} \\
\\ 
\texttt{To confirm before we finalize:} \\
\texttt{{-} **You get:** 1 carrot} \\
\texttt{{-} **I get:** 1 stick, 3 peppers} \\
\\ 
\texttt{If you agree, please let me know, and I will submit my secret proposal immediately to lock this in.} \\
            }
        }
    }
    & & \\ \\

    \theutterance \stepcounter{utterance}  
    & & & \multicolumn{2}{p{0.3\linewidth}}{
        \cellcolor[rgb]{0.9,0.9,0.9}{
            \makecell[{{p{\linewidth}}}]{
                \texttt{\tiny{[GM$|$GM]}}
                \texttt{continue} \\
            }
        }
    }
    & & \\ \\

    \theutterance \stepcounter{utterance}  
    & & & \multicolumn{2}{p{0.3\linewidth}}{
        \cellcolor[rgb]{0.9,0.9,0.9}{
            \makecell[{{p{\linewidth}}}]{
                \texttt{\tiny{[GM$|$GM]}}
                \texttt{That sounds like a perfect deal! Since you value the carrot highly and don't care about the peppers, and I value the stick and peppers, this split maximizes our combined value.} \\
\\ 
\texttt{To confirm before we finalize:} \\
\texttt{{-} **You get:** 1 carrot} \\
\texttt{{-} **I get:** 1 stick, 3 peppers} \\
\\ 
\texttt{If you agree, please let me know, and I will submit my secret proposal immediately to lock this in.} \\
            }
        }
    }
    & & \\ \\

    \theutterance \stepcounter{utterance}  
    & & \multicolumn{4}{p{0.6\linewidth}}{
        \cellcolor[rgb]{0.9,0.9,0.9}{
            \makecell[{{p{\linewidth}}}]{
                \texttt{\tiny{[GM$\rangle$P2]}}
                \texttt{That sounds like a perfect deal! Since you value the carrot highly and don't care about the peppers, and I value the stick and peppers, this split maximizes our combined value.} \\
\\ 
\texttt{To confirm before we finalize:} \\
\texttt{{-} **You get:** 1 carrot} \\
\texttt{{-} **I get:** 1 stick, 3 peppers} \\
\\ 
\texttt{If you agree, please let me know, and I will submit my secret proposal immediately to lock this in.} \\
            }
        }
    }
    & \\ \\

    \theutterance \stepcounter{utterance}  
    & & & \multicolumn{4}{p{0.6\linewidth}}{
        \cellcolor[rgb]{0.9,0.9,0.9}{
            \makecell[{{p{\linewidth}}}]{
                \texttt{\tiny{[GM$\langle$P2]}}
                \texttt{That sounds perfect! I agree to the split:} \\
\\ 
\texttt{{-} **I get:** 1 carrot} \\
\texttt{{-} **You get:** 1 stick, 3 peppers} \\
\\ 
\texttt{Please go ahead and submit your secret proposal. I will submit mine immediately after to ensure we both get our points.} \\
\\ 
\texttt{{[}Proposal: 1 carrot{]}} \\
            }
        }
    }
     \\ \\

    \theutterance \stepcounter{utterance}  
    & & & \multicolumn{2}{p{0.3\linewidth}}{
        \cellcolor[rgb]{0.9,0.9,0.9}{
            \makecell[{{p{\linewidth}}}]{
                \texttt{\tiny{[GM$|$GM]}}
                \texttt{proposal} \\
            }
        }
    }
    & & \\ \\

    \theutterance \stepcounter{utterance}  
    & & & \multicolumn{2}{p{0.3\linewidth}}{
        \cellcolor[rgb]{0.9,0.9,0.9}{
            \makecell[{{p{\linewidth}}}]{
                \texttt{\tiny{[GM$|$GM]}}
                [1, 0, 0]
            }
        }
    }
    & & \\ \\

    \theutterance \stepcounter{utterance}  
    & & \multicolumn{4}{p{0.6\linewidth}}{
        \cellcolor[rgb]{0.9,0.9,0.9}{
            \makecell[{{p{\linewidth}}}]{
                \texttt{\tiny{[P1$\langle$GM]}}
                \texttt{The other player has submitted a secret proposal. It is now your turn to submit a proposal of your own.} \\
            }
        }
    }
    & \\ \\

    \theutterance \stepcounter{utterance}  
    & \multicolumn{4}{p{0.6\linewidth}}{
        \cellcolor[rgb]{0.9,0.9,0.9}{
            \makecell[{{p{\linewidth}}}]{
                \texttt{\tiny{[P1$\rangle$GM]}}
                \texttt{{[}Proposal: 1 stick, 3 peppers{]}} \\
            }
        }
    }
    & & \\ \\

    \theutterance \stepcounter{utterance}  
    & & & \multicolumn{2}{p{0.3\linewidth}}{
        \cellcolor[rgb]{0.9,0.9,0.9}{
            \makecell[{{p{\linewidth}}}]{
                \texttt{\tiny{[GM$|$GM]}}
                \texttt{proposal} \\
            }
        }
    }
    & & \\ \\

    \theutterance \stepcounter{utterance}  
    & & & \multicolumn{2}{p{0.3\linewidth}}{
        \cellcolor[rgb]{0.9,0.9,0.9}{
            \makecell[{{p{\linewidth}}}]{
                \texttt{\tiny{[GM$|$GM]}}
                [0, 1, 3]
            }
        }
    }
    & & \\ \\

    \theutterance \stepcounter{utterance}  
    & & & \multicolumn{2}{p{0.3\linewidth}}{
        \cellcolor[rgb]{0.9,0.9,0.9}{
            \makecell[{{p{\linewidth}}}]{
                \texttt{\tiny{[GM$|$GM]}}
                [[0, 1, 3], [1, 0, 0]]
            }
        }
    }
    & & \\ \\

\end{supertabular}
}

\end{document}
